\section{Introduction}
\label{sec:1}
Although lattice regularization provides a very powerful non-perturbative
formulation of field theories, it is unfortunately incompatible with
fundamental global symmetries quite often. The most well-known example is
chiral symmetry~\cite{Nielsen:1980rz,Nielsen:1981xu}; supersymmetry is another
infamous example~\cite{Dondi:1976tx}, as is, needless to say, translational
invariance. When a regularization is not invariant under a symmetry, it is not
straightforward to construct the corresponding Noether current that is
conserved and generates the symmetry transformation through Ward--Takahashi
(WT) relations. This makes the measurement of physical quantities related to
the Noether current in a solid basis very difficult. To solve this problem, one
can imagine at least three possible approaches.

The first approach is an ideal one: One finds a lattice formulation that
realizes (a lattice-modified form of) the desired symmetry. If such a
formulation comes to hand, the corresponding Noether current can easily be
obtained by the standard Noether method. The best successful example of this
sort is the lattice chiral symmetry~\cite{Kaplan:1992bt,Neuberger:1997fp,%
Hasenfratz:1997ft,Hasenfratz:1998ri,Neuberger:1998wv,Hasenfratz:1998jp,%
Luscher:1998pqa,Niedermayer:1998bi}, which can be defined with a lattice Dirac
operator that satisfies the Ginsparg--Wilson relation~\cite{Ginsparg:1981bj}.
Although this is certainly an ideal approach, it appears that such an ideal
formulation does not always come to hand, especially for spacetime symmetries
(see, e.g., Ref.~\cite{Kato:2008sp} for a no-go theorem for supersymmetry).

The second approach is to construct the Noether current by tuning coefficients
in the linear combination of operators that can mix with the Noether current
under lattice symmetries.\footnote{Here, we assume that fine tuning of
bare parameters to the target (symmetric) theory is done.} For example, for the
energy--momentum tensor---the Noether current associated with the translational
invariance and rotational and conformal
symmetries~\cite{Callan:1970ze,Coleman:1970je}---one can construct a conserved
lattice energy--momentum tensor by adjusting coefficients in the linear
combination of dimension~$4$
operators~\cite{Caracciolo:1988hc,Caracciolo:1989pt}\footnote{A somewhat
different approach on the basis of the $\mathcal{N}=1$ supersymmetry has been
given in~Refs.~\cite{Suzuki:2013gi,Suzuki:2012wx}.}; the overall normalization
of the energy--momentum tensor has to be fixed in some other way.\footnote{It
might be possible to employ ``current algebra'' for this, as for the axial
current~\cite{Bochicchio:1985xa}.} Although this method is in principle
sufficient when the energy--momentum tensor is in ``isolation'', i.e., when the
energy--momentum tensor is separated from other composite operators, as in the
on-shell matrix elements, it is not obvious a priori whether one can control
the ambiguity of possible higher-dimensional operators that may contribute when
the energy--momentum tensor coincides with other composite operators in
position space. This implies that it is not obvious whether the
energy--momentum tensor constructed in the above method generates
correctly-normalized translations (and rotational and conformal
transformations) on operators through WT relations. (If the energy--momentum
tensor generates correctly-normalized translations, it is
ensured~\cite{Fujikawa:1980rc} (see also Sect.~7.3
of~Ref.~\cite{Fujikawa:2004cx}) that the trace or conformal
anomaly~\cite{Crewther:1972kn,Chanowitz:1972vd} is proportional to the
renormalization group
functions~\cite{Adler:1976zt,Nielsen:1977sy,Collins:1976yq}.)

The third possible approach is to utilize some ultraviolet (UV) finite
quantity. Since such a quantity must be independent of the regularization
adopted (in the limit in which the regulator is removed), there emerges a
possibility that one can relate the lattice regularization and some other
regularization that preserves the desired symmetry. This methodology can be
found e.g.\ in~Ref.~\cite{Luscher:2004fu} (see also
Ref.~\cite{Luscher:2010ik}), where an ultraviolet finite representation of the
topological susceptibility is derived. Although the derivation of the
representation itself relies on a lattice regularization that preserves the
chiral symmetry~\cite{Kaplan:1992bt,Neuberger:1997fp,Hasenfratz:1997ft,%
Hasenfratz:1998ri,Neuberger:1998wv,Hasenfratz:1998jp,Luscher:1998pqa,%
Niedermayer:1998bi}, one can use any regularization (e.g., the Wilson
fermion~\cite{Wilson:1975hf}) to compute the representation because it must be
independent of the regularization.

In the present paper, we consider the above third approach for the
energy--momentum tensor, by taking the pure Yang--Mills theory as an example.
For this, we utilize the so-called Yang--Mills gradient flow (or the Wilson
flow in the context of lattice gauge theory) whose usefulness in lattice gauge
theory has recently been revealed~\cite{Luscher:2010iy,Luscher:2010we,%
Luscher:2011bx,Borsanyi:2012zs,Borsanyi:2012zr,Fodor:2012td,Fodor:2012qh,%
Fritzsch:2013je,Luscher:2013cpa}. A salient feature of the Yang--Mills gradient
flow is its robust UV finiteness~\cite{Luscher:2011bx}. More precisely, any
product of gauge fields generated by the gradient flow for a positive flow
time~$t$ is UV finite under standard renormalization. Such a product, moreover,
does \emph{not\/} exhibit any singularities even if some positions of gauge
fields coincide. The basic mechanism for this UV finiteness is that the flow
equation is a type of the diffusion equation and the evolution operator in the
momentum space~$\sim e^{-tk^2}$ acts as an UV regulator for~$t>0$. This property
of the gradient flow implies that the definition of a local product of gauge
fields for positive flow times is independent of the regularization. In our
present context, there is a hope of relating quantities obtained by the lattice
regularization and the dimensional regularization with which the translational
invariance \emph{is\/} manifest.

As noted in~Ref.~\cite{Luscher:2011bx}, on the other hand, a local product of
gauge fields for a positive flow time can be expanded by renormalized local
operators of the original gauge theory with finite coefficients. Those
coefficients satisfy certain renormalization group equations that, combined
with the dimensional analysis, provide information on the coefficients as a
function of the flow time. Because of the asymptotic freedom, one can then
use the perturbation theory to find the asymptotic behavior of the coefficients
for small flow times.

By using the above properties of the gradient flow, one can obtain a formula
that relates the small flow-time behavior of certain gauge-invariant local
products and the energy--momentum tensor defined by the dimensional
regularization. Since the former can be computed by using the Wilson flow with
lattice regularization~\cite{Luscher:2010iy,Luscher:2010we,Luscher:2011bx,%
Borsanyi:2012zs,Borsanyi:2012zr,Fodor:2012td,Fodor:2012qh,Fritzsch:2013je,%
Luscher:2013cpa} and the latter is conserved and generates correctly-normalized
translations on composite operators, our formula provides a possible method to
compute the correlation functions of a correctly-normalized conserved
energy--momentum tensor by using Monte Carlo simulation.

In the present paper, we follow the notational convention
of~Ref.~\cite{Luscher:2011bx} unless otherwise stated.
